\section{Results}
\suppressfloats[t]
\label{sec:results}

\subsection{RQ1: Comparative information expands the finite-policy frontier only at a relaxed tolerance}
\label{sec:rq1-answerability}

\textbf{Answer.} On the target-fixed held-out grid, the target-only forced
policy answered every scope arm but had \VFiveTargetForcedRiskPercent{} observed
conditional scope error. Thus no predeclared target-only policy had positive
finite-policy coverage at any reported tolerance below $0.5$. The comparative
forced policy also answered every arm but had
\VFiveComparativeForcedRiskPercent{} observed error. Its confidence-selective
policies calibrated at $.01$, $.05$, and $.10$ answered
\VFiveStrictConfidenceAnsweredArms{} held-out arms. Only the policy calibrated
at $.20$ reached a positive comparative envelope:
\VFiveComparativeFrontierTwentyPercent{} coverage at
\VFiveComparativeRiskTwentyPercent{} observed conditional error.

\begin{figure}[tbp]
\centering
\includegraphics[width=0.88\linewidth]{fig_v05_answerability_frontier.png}
\caption{Held-out finite-policy answerability frontiers by information
channel. Target-only policies do not reach positive coverage at the four
reported tolerances. Comparative coverage emerges only at $\alpha=.20$.
The certificate-assisted policy is plotted separately because it uses bounded
synthetic-design information, not a real-network deployment channel.}
\label{fig:v05-answerability-frontier}
\end{figure}

\begin{table}[tbp]
\centering
\caption{Frozen finite-policy scope-answerability envelope on the
held-out v0.5 evaluation components. Parentheses give observed
conditional scope error for the policy attaining the displayed coverage;
\texttt{--} means no positive-coverage candidate qualified. The
certificate-assisted channel uses synthetic design information and is not
an operational deployment channel. Its 0 observed errors are counted
among 179,994 answered scope arms, not estimated
population risk.}
\label{tab:v05-frontier}
\small
\begin{tabular}{lrrr}
\toprule
Tolerance $\alpha$ & Target-only & Comparative & Certificate-assisted \\
\midrule
0.01 & 0\% & 0\% (--) & 39.1\% (0 observed errors) \\
0.05 & 0\% & 0\% (--) & 39.1\% (0 observed errors) \\
0.10 & 0\% & 0\% (--) & 39.1\% (0 observed errors) \\
0.20 & 0\% & 60.4\% (19.5\%) & 39.1\% (0 observed errors) \\
\bottomrule
\end{tabular}
\end{table}


At $\alpha=.20$, the achieved finite-policy Scope Answerability Gain from the
target-only to comparative channel is
\VFiveComparativeFrontierTwentyPercent{}. Its component-bootstrap interval is
$[\VFiveComparativeGainLowerPercent{},\VFiveComparativeGainUpperPercent{}]$.
This is a comparison among frozen policies on held-out synthetic components,
not an estimate of a universal optimum or future population risk.

\FloatBarrier

\subsection{RQ2: Structural separation maps a retained abstention boundary}
\label{sec:rq2-certificate}

\textbf{Answer.} The structural interval-separation certificate answered
\VFiveCertificateAnsweredPairs{} of \VFiveEvaluationPairs{} matched pairs
(\VFiveCertificateCoveragePercent{}), with
\VFiveCertificateObservedErrorDisplay{} among
\VFiveCertificateAnsweredArms{} answered scope arms. It recovered
\VFiveCertificateEfficiencyPercent{} of pairs in the predeclared
simulation-information reference region and abstained on the remaining
\VFiveNonpositiveMarginPairs{} pairs (\VFiveCertificateAbstentionPercent{}),
including all \VFiveQZeroPairs{} observationally indistinguishable $q=0$
negative controls. It had \VFiveEnvelopeViolations{} observed envelope
violations. These finite-grid observations hold only for the bounded generator:
they neither establish zero population risk nor yield a deployable
real-network certificate.

\begin{figure}[tbp]
\centering
\includegraphics[width=\linewidth]{fig_v05_scope_boundary.png}
\caption{Structural answerability boundary on the held-out grid. From left to
right, the panels show the median normalized structural margin, certificate
answered-pair coverage, and forced-comparative scope error. Values above one in
the first panel meet the sufficient interval-separation condition. Each cell
averages over the five adverse comparison factors; $q=0$ remains the
observational-identity negative control.}
\label{fig:v05-scope-boundary}
\end{figure}

\begin{table}[tbp]
\centering
\caption{Certificate behavior by nominal donor participation on held-out
v0.5 pairs. Efficiency is coverage relative to the predeclared
simulation-information oracle region; it is undefined at $q=0$ because
that negative-control stratum has no oracle-answerable pairs.}
\label{tab:v05-certificate}
\small
\begin{tabular}{lrrr}
\toprule
Participation & Pair coverage & Conditional error & Oracle efficiency \\
\midrule
$q=0.00$ & 0.000000 & -- & -- \\
$q=0.25$ & 0.109375 & 0.000000 & 0.875000 \\
$q=0.50$ & 0.382812 & 0.000000 & 0.942308 \\
$q=0.75$ & 0.663997 & 0.000000 & 0.988435 \\
$q=1.00$ & 0.796875 & 0.000000 & 1.000000 \\
\bottomrule
\end{tabular}
\end{table}


The certificate coverage increases with nominal donor participation and signal,
but it is intentionally incomplete under contamination, nonlinear raw fields,
and heteroskedastic noise. In particular, the $q=0$ stratum contains
\VFiveQZeroPairs{} matched pairs with identical target and donor observations
in both arms, produces \VFiveQZeroCertificateAnsweredPairs{} certificate
answers, and has no reference-region pair authorized for an answer. The result
does not declare all real partial scope objectively ambiguous; it marks where
this sufficient bound does not authorize an answer.

\begin{table}[tbp]
\centering
\caption{Predeclared failure and abstention accounting on
230,400 held-out v0.5 matched pairs. These rows are retained
rather than converted into a favorable single score.}
\label{tab:v05-failure-accounting}
\small
\begin{tabularx}{\linewidth}{@{}p{0.42\linewidth}rX@{}}
\toprule
Outcome & Count & Interpretation \\
\midrule
Target and donor observations identical at $q=0$ & 46,080 & negative-control abstention stratum \\
Nonpositive structural margin & 140,403 & certificate must abstain \\
Certificate-answered pairs & 89,997 & 0 errors among 179,994 answered arms \\
Envelope violations & 0 & none observed \\
Comparative-forced errors & 126,764 & among 460,800 forced scope arms \\
Confidence policy at calibration $\alpha=0.05$ & 0 & no answered arms \\
\bottomrule
\end{tabularx}
\end{table}


\FloatBarrier

\subsection{RQ3: Cross-site comparison separates the constructed task, but estimator superiority is not supported}
\label{sec:rq3-methods}

\textbf{Answer.} In the independent stable-regime benchmark, all four
cross-site methods have aggregate AUPRC between \CrossSiteAUPRCLow{} and
\CrossSiteAUPRCHigh{}, whereas the five target-only methods range only from
\TargetOnlyAUPRCLow{} to \TargetOnlyAUPRCHigh{}
(Table~\ref{tab:all-methods} in Appendix~\ref{sec:appendix-methods}). This
positive information contrast is the intended behavior of the constructed
local-versus-regional task, not a generic change-detection claim.

\begin{figure}[tbp]
\centering
\includegraphics[width=\linewidth]{fig_cross_site_scope_metrics.pdf}
\caption{Cross-site versus target-only ranking on the independent held-out
stable-regime benchmark. All four cross-site methods separate the constructed
local-versus-shared scope task, while the five target-only methods remain near
chance. The comparison demonstrates an information contrast in this benchmark,
not a universal method ranking or a real-network mechanism result.}
\label{fig:cross-site-scope-metrics}
\end{figure}

\begin{table}[!ht]
\centering
\caption{Independent stable-regime synthetic evaluation.}
\label{tab:synthetic-metrics}
\small
\begin{tabular}{lrrrr}
\toprule
Method & MAE & AUPRC & Macro-F1 & Regional FPR \\
\midrule
Standard SC & 0.100 & 0.915 & 0.816 & 0.140 \\
MetaShift fixed & 0.103 & 0.898 & 0.795 & 0.150 \\
MetaShift CV & 0.097 & 0.902 & 0.809 & 0.135 \\
Nearest-neighbor DiD & 0.111 & 0.884 & 0.789 & 0.075 \\
\bottomrule
\end{tabular}
\par\smallskip
\begin{minipage}{0.94\linewidth}
\footnotesize\textit{Note.} Metrics are aggregates over the frozen \texttt{stable\_full\_v2} evaluation set. Lower MAE and regional false-positive rate are preferable; higher AUPRC and macro-F1 are preferable.
\end{minipage}
\end{table}


Standard synthetic control has the strongest aggregate AUPRC and macro-F1.
Cross-validated MetaShift has the smallest aggregate local-effect MAE, while
nearest-neighbor difference-in-differences has the lowest regional false-positive
rate. No method wins every metric. Most importantly, neither MetaShift variant
has a confidence-supported aggregate MAE advantage over standard synthetic
control: both paired component-bootstrap intervals include zero
(Table~\ref{tab:paired-bootstrap}).

\begin{table}[tbp]
\centering
\caption{Paired event-cluster bootstrap for local-effect MAE differences.}
\label{tab:paired-bootstrap}
\small
\begin{tabular}{lrrr}
\toprule
Comparison & Difference & 95\% bootstrap interval & Clusters \\
\midrule
MetaShift fixed minus Standard SC & 0.00361 & [-0.00401, 0.01163] & 80 \\
MetaShift CV minus Standard SC & -0.00240 & [-0.00772, 0.00347] & 80 \\
\bottomrule
\end{tabular}
\par\smallskip
\begin{minipage}{0.94\linewidth}
\footnotesize\textit{Note.} Difference is MetaShift minus standard synthetic control; negative values favor MetaShift on MAE. Both bootstrap intervals include zero.
\end{minipage}
\end{table}


\begin{figure}[tbp]
\centering
\includegraphics[width=0.88\linewidth]{fig_paired_bootstrap.pdf}
\caption{Paired event-cluster bootstrap intervals for MetaShift minus standard
synthetic-control local-effect MAE. Negative values favor MetaShift; both
intervals cross zero. The frozen evidence therefore does not support an
aggregate MetaShift-superiority claim.}
\label{fig:paired-bootstrap}
\end{figure}

\FloatBarrier

\subsection{RQ4: A qualified AQS comparison is available for 228 of 563 anchors}
\label{sec:rq4-deployment}

\textbf{Answer.} The distinct-physical-donor rule permits a complete common
cross-site comparison for \CompleteComparisons{} of \TotalAnchors{} anchors
(40.5\%). The other \InsufficientDonorAnchors{} donor-insufficient and
\InputFailureAnchors{} input-window failures are retained as abstentions, not
dropped observations. This availability result is the AQS audit's primary
deployment finding.

\begin{table}[!ht]
\centering
\caption{Complete 88101 event audit and real-event diagnostics.}
\label{tab:real-audit}
\small
\begin{tabular}{lr}
\toprule
Audit disposition & Anchors \\
\midrule
Complete common-method comparison & 228 \\
Fewer than three distinct physical donors & 325 \\
Estimator input-window failure & 10 \\
Total & 563 \\
\bottomrule
\end{tabular}
\par\medskip
\begin{tabular}{lrrr}
\toprule
Method & Events & Fixed conditional interval excludes 0 & Median log effect \\
\midrule
MetaShift fixed & 228 & 159 & -0.087 \\
Standard SC & 228 & 145 & -0.071 \\
Nearest-neighbor DiD & 228 & 134 & -0.075 \\
\bottomrule
\end{tabular}
\par\smallskip
\begin{minipage}{0.94\linewidth}
\footnotesize\textit{Note.} Nested selection-aware intervals are available for 227/228 complete events; 153 exclude zero. Their mean width is 0.191 log units versus 0.179 for fixed-weight MetaShift intervals. Leave-one-donor-out refits complete for 227 events, with 202 retaining direction under every donor removal. All intervals are diagnostic, not calibrated physical-bias confidence intervals.
\end{minipage}
\end{table}


\begin{figure}[tbp]
\centering
\includegraphics[width=0.88\linewidth]{fig_event_accounting.pdf}
\caption{AQS audit accounting. Panel (a) partitions all \TotalAnchors{}
metadata anchors into donor-insufficient, input-window-failure, and complete
comparison dispositions. Panel (b) partitions complete comparisons into
protocol evidence tiers. These are audit dispositions, not diagnoses of
physical monitor changes.}
\label{fig:event-accounting}
\end{figure}

The exploratory evidence synthesis assigns \SupportedCandidates{} supported
candidates, \NotSupportedEvents{} not-supported events, and
\InconclusiveEvents{} inconclusive events. ``Supported'' means that all
predeclared diagnostic conditions were available and passed; it does not
confirm a physical failure, replacement, measurement bias, or pollution
correction.

\FloatBarrier

\subsection{RQ5: Calibration and contextual checks limit, rather than validate, real-data interpretation}
\label{sec:rq5-limits}

\textbf{Answer.} The fixed-protocol interval constructions do not support a
generic calibration claim. On the held-out synthetic effects, nominal 95\%
fixed-weight coverage ranges from 63.9\% to 67.3\% over
\CoverageInstancesPerMethod{} instances per method, while nominal 90\%
split-conformal coverage ranges from 98.8\% to 99.6\% and is conservative.

\begin{figure}[tbp]
\centering
\includegraphics[width=\linewidth]{fig_interval_coverage.pdf}
\caption{Held-out synthetic interval coverage and mean width. Fixed conditional
intervals under-cover their nominal 95\% target, whereas split-conformal
intervals over-cover their nominal 90\% target. The figure reports a
calibration limitation, not a real-event confidence guarantee.}
\label{fig:interval-coverage}
\end{figure}

The real-data date-resampling diagnostic holds each event's pre-transition
donor weights fixed and compares the mean standardized score at the metadata
anchor with one randomly selected stable post-transition date per eligible
event. Under the within-event exchangeability diagnostic null, its statistic is
the difference between the anchor mean and sampled-placebo mean; the reported
one-sided tail probability is
$(1+\#\{\text{sampled placebo mean}\geq\text{anchor mean}\})/201 =
\DateResamplingP{}$. It is a comparison with stable dates, not a causal
probability or a physical-mechanism test.

\begin{table}[tbp]
\centering
\caption{Placebo, external-document, and same-site POC evidence boundaries.}
\label{tab:placebo-external}
\scriptsize
\begin{tabular}{lr}
\toprule
Diagnostic & Value \\
\midrule
Complete time-placebo events with at least 50 dates & 157 \\
Complete time-placebo events with 100 dates & 128 \\
Raw time-placebo probabilities at most 0.10 & 71 \\
BH q values at most 0.10 & 40 \\
Donor-as-treated placebo records & 866 \\
Date-resampling permutations & 200 \\
Dated site-specific external confirmations & 0/20 \\
Usable hourly same-site POC comparisons & 9/11 \\
Daily/hourly POC direction agreements & 8/9 \\
Standard SC 90th-percentile risk-gate retention & 75/80 \\
\bottomrule
\end{tabular}
\par\smallskip
\begin{minipage}{0.94\linewidth}
\footnotesize\textit{Note.} Time and donor placebos are diagnostic falsifications. Same-site POC and document review provide contextual evidence only, not instrument ground truth.
\end{minipage}
\end{table}


Of \POCCandidates{} same-site POC candidates, \POCUsablePairs{} have usable
paired hourly windows and \POCDirectionAgreements{} agree in daily/hourly
direction. Targeted official-document review finds
\ExternalReviewConfirmations{}/\ExternalReviewEvents{} dated site-specific
confirmations. Missing documentation and unavailable POC data remain
unavailable evidence; they are not evidence that nothing changed.

\FloatBarrier
